\documentclass{article}
\usepackage{amsmath,amssymb}
\usepackage[vocab={kinetics}]{leanify}
\newtheorem{lemma}{Lemma}
\begin{document}

\section{Notation}

We write $Z(t,k,c)$ for the concentration at time $t$ in a zeroth-order reaction with rate constant $k$ and initial concentration $c$.
We write $F(t,k,c)$ for the concentration at time $t$ in a first-order reaction with rate constant $k$ and initial concentration $c$.
We write $S(t,k,c)$ for the concentration at time $t$ in a second-order reaction with rate constant $k$ and initial concentration $c$.
We write $H(k)$ for the half-life of a first-order reaction with rate constant $k$.
We write $J(k,c)$ for the half-life of a second-order reaction with rate constant $k$ and initial concentration $c$.
We write $\lambda(T)$ for the decay constant for the half-life $T$.
We write $A(a,E,T,R)$ for the Arrhenius rate constant for the prefactor $a$, the activation energy $E$ and the temperature $T$ with gas constant $R$.
We write $v(V,K,s)$ for the Michaelis-Menten rate for the maximal rate $V$, the Michaelis constant $K$ and the substrate concentration $s$.
We write $P(x)$ for the pH for the hydrogen ion concentration $x$.

\section{Half-lives}

\begin{lemma}\label{half}
Let $k$ and $c$ be positive real numbers. Then $F(H(k),k,c)=\frac{c}{2}$, $S(J(k,c),k,c)=\frac{c}{2}$ and $Z(\frac{c}{k},k,c)=0$.
\end{lemma}

\begin{lemma}\label{decay}
Let $T$ be a positive real number. Then $H(\lambda(T))=T$.
\end{lemma}

\section{Rates}

\begin{lemma}\label{rates}
$P(\frac{1}{1000})=3$, $v(2,1,1)=1$ and $A(1,0,1,1)=1$.
\end{lemma}

\end{document}
